\section{准 TMD 因子化的 Collins 式证明与匹配核}\label{sec:yzliu-quasi-tmd}

准 TMD 因子化定理陈述于第~22~章（\cref{thm:lamet-quasi-tmd}，[A]）。本节给出 Liu 的 Collins 式证明路线——单线记号下的幂次计数（\cref{fig:yzliu-leading-region}）——以及硬核的全阶结构与显式结果。

%%FIG fig:yzliu-leading-region

\begin{proposition}[准 TMD 因子化的 Collins 式证明：单线记号下的幂次计数 \tagA]\label{prop:yzliu-qtmd-powercount}
在以在壳无质量部分子为外态的微扰论各阶，对\cref{def:lamet-quasi-tmd}的准 TMD，在局部假设 (QT1)--(QT7) 之下：(i) Landau 方程给出的领头区域只有一个外共线方向，共线位移与订书钉位移退耦，订书钉传播子只能经偶极子 $2\to2$ 插入（\cref{prop:yzliu-pinch-ultrasoft}的捏点）取得非零的 Feynman 参数解；(ii) 在单线记号（一对订书钉线 = 一条量纲 $\frac32$、自旋 $0$ 的线）下，幂次计数参数 $\lambda_{\mathrm{pc}}\equiv1/(b_\perp P^z)$ 的幂为 $\lambda_{\mathrm{pc}}^{F(H)+F(SC)+F(H')+F(H'')}$，% VF-YZL-8
\begin{equation}\label{eq:yzliu-FH}
\begin{aligned}
F(H)&=\sum_{s=0,1}\big[N(CH,1,s)+N(SH,1,s)\big]+N(CH,\tfrac32,\tfrac12)\\
&\quad+\tfrac32\sum_{s=0,\frac12}N(SH,\tfrac32,s)-N(CH,1,\ell)-\tfrac72 ,\qquad F(SC)=\tfrac12N(SC,\tfrac32,\tfrac12)+N(SC,1,0),
\end{aligned}
\end{equation}
$F(H')=F(H'')=0$（$N(XY,d,s)$ 为子图 $X$、$Y$ 之间量纲 $d$、自旋 $s$ 的线数，$\ell$ 指纵向极化的共线胶子）；$F(H),F(SC)\ge0$（无超领头区域），领头区域为 $N(CH,\frac32,\frac12)=2$、$N(SH,\frac32,0)=1$、任意多条纵向共线胶子进入 $H$、任意多条软胶子进入 $C$，其余为零（$F(H)=-\frac72+2+\frac32=0$；每多一条进入 $H$ 的软胶子代价 $+1$）；(iii) 经软—共线（Grammer--Yennie 型）近似与 Ward 恒等式，在任何光锥上的正规子中 $\tilde f=H\,\frac{f(x,b_\perp,\mu,\delta^-)}{S(b_\perp,\mu,\delta^+,\delta^-)}\,S(b_\perp,\mu,\delta^+,\hat z)$\cite[式~(74)]{Liu:2022nnk}；(iv) 与\cref{thm:yzliu-intrinsic-soft}结合即得\cref{thm:lamet-quasi-tmd}（指数为 $\frac12\tilde K\ln\frac{\zeta_z}{\zeta}$），动量演化 $P^z\frac{\dd}{\dd P^z}\ln\tilde f=\tilde K+\mathcal G$，$\mathcal G^{(1)}=\frac{\alpha_sC_F}{\pi}(1-L_\zeta)$。% VF-YZL-5
局部假设（原文的 (Q1)--(Q7)）：(QT1) 领头区域由 Landau 方程给出；(QT2) 单个割图中可能出现的其他共线方向在对割求和后相消；(QT3) 无 Glauber 区：硬区内所有小 $k^z$ 的规范线极点可变形到上半平面，同一变形避开 $k^+k^-\ll k_\perp^2$；(QT4) 捏点与超软模式如\cref{prop:yzliu-pinch-ultrasoft}；(QT5) 硬子图是均匀的，各分量 $\sim P^z$；(QT6) 订书钉算符乘法可重整（\cref{asm:lamet-staple}(ST1)）；(QT7) 以在壳无质量部分子为外态的微扰论。（本命题只涉及部分子态。）
% errtype: E-S
\uses{def:lamet-quasi-tmd,thm:lamet-quasi-tmd,asm:lamet-staple,prop:yzliu-pinch-ultrasoft,thm:yzliu-intrinsic-soft,prop:tmd-regions}
\end{proposition}

\begin{proofsketch}
（1）Landau 方程\cite[附录~C]{Liu:2022nnk}与对割求和（(QT2)，与\cref{prop:tmd-regions}同型）给出 (i)。（2）$F(H)=-D(H)-B(H)+1$：硬核的质量量纲 $D(H)=\frac92-\sum_iN_id_i$（$\frac92$ 来自两个夸克场与一条单线订书钉），推动增强 $B(H)=N(CH,1,\ell)+\frac12N(CH,\frac32,\frac12)$，$+1$ 来自归一化 $1/P^z$，即得\cref{eq:yzliu-FH}；孤立硬核 $H'$、$H''$ 可吸收进软区、共线区或重整化\cite[式~(54)--(61)]{Liu:2022nnk}。（3）进入 $H$ 的总有两条夸克线与一条单线订书钉，且 $N(CH,1,1)\ge N(CH,1,\ell)$，故 $F(H)\ge0$，等号即领头区域；硬动量不能在 $0$ 与 $\vec b_\perp$ 之间传递，故没有 $x$ 卷积\cite[式~(64)--(68)]{Liu:2022nnk}（此步只给出 E-S 型陈述）。（4）硬区内规范线传播子为 $\prod_j(\sum_{k\ge j}p^z_k+\ii0)^{-1}$，$p^z_i=\sum_ja_{ij}k^z_j$，$a_{ij}\ge0$，故所有 $k^z_j$ 可变形到上半平面，同时给出 (QT5) 与 (QT3)\cite[式~(62)--(63)]{Liu:2022nnk}。（5）软—共线近似与 Ward 恒等式给出 (iii)。（6）由\cref{thm:yzliu-intrinsic-soft}(b)，(iii) 中的正规子与 $f^{\mathrm{TMD}}=f(\delta^-)/\sqrt{S(\delta^+,\delta^-)}$ 的正规子相消，常数项给出 $1/\sqrt{S_I}$，$\zeta_z$ 依赖由 CS 方程与动量演化固定；单圈检验见\cite[式~(23)、(39)、(85)]{Liu:2022nnk}。
\end{proofsketch}

\paragraph{地位与缺口。}全阶论证 [A]：幂次计数各步只给出 E-S 型陈述，(QT1)--(QT5) 是 CSS 型论证的标准步骤而非定理；PLB 版本自称是数学证明的“sketch”\cite{Ji:2019ewn}；提升为 [R] 需要\cref{sec:yzliu-intro}所引的多标度分析。据我们截至 2026 年 10 月的文献检索（INSPIRE 检索），尚未见针对有限长订书钉的专门 Glauber 分析。独立的支持：Vladimirov 与 Schäfer 的 SCET$_{\mathrm{II}}$ 推导（附单圈检验）\cite{Vladimirov:2020ofp}，Ebert、Schindler、Stewart 与 Zhao 经大快度方案的全阶论证\cite{Ebert:2022fmh}，NNLO 的直接计算\cite{delRio:2023pse}与\cite{Ji:2023pba}。Ji、Liu 与 Liu\cite{Ji:2019ewn}把准 TMD 匹配到方案无关、软因子可欧氏计算的 TMD 分布；单线幂次计数见\cite[式~(54)--(68)]{Liu:2022nnk}。

\begin{proposition}[准 TMD 匹配核：RGE、重—轻分解与两圈结果 \tagA]\label{prop:yzliu-hard-kernel}
(a) 硬核满足
\begin{equation}\label{eq:yzliu-H-rge}
\mu\frac{\dd}{\dd\mu}\ln H\Big(\frac{\zeta_z}{\mu^2}\Big)=\Gamma_{\mathrm{cusp}}\ln\frac{\zeta_z}{\mu^2}+\gamma_C,\qquad \gamma_C\equiv\tilde\gamma_H=\gamma_V+2\gamma_F+\gamma_{HH}=\gamma_V+2\gamma_F+2\gamma_{HL}-\gamma_s ,
\end{equation}
$\gamma_V$ 为轻夸克 Sudakov 形状因子的非尖点反常维度，$\gamma_F$ 为重—轻流的反常维度，$\gamma_{HL}$ 为重—轻尖点的常数，$\gamma_s$ 为类光软常数\cite[式~(A.11)]{Ji:2023pba}（PLB 的 $2\gamma_F+\Gamma_S+2\gamma_H$ 与讲义的 $2\gamma_F-\Gamma_S+2\gamma_J$ 按\cref{tab:yzliu-conventions}换算后与此一致）。(b)（单圈内容，\tagR，$R_{\mathrm{PT}}$）（$R_{\mathrm{PT}}$：在重整化微扰论中对红外正规化的部分子态严格。）% VF-YZL-6
\begin{equation}\label{eq:yzliu-H1}
H=1+a\Big(-2+\frac{\pi^2}{12}-\frac{L_\zeta^2}{2}+L_\zeta\Big)=1+\frac{\alpha_sC_F}{4\pi}\Big(-L_\zeta^2+2L_\zeta-4+\frac{\pi^2}{6}\Big),\qquad H=C^+_{HL}C^-_{HL},
\end{equation}
$C^\pm_{HL}=1+\frac{\alpha_sC_F}{8\pi}(-4-\frac{5\pi^2}{6}+2L_\pm-L_\pm^2)$，$L_\pm=L_\zeta\pm\ii\pi$，$\Imag C^{\pm(1)}_{HL}=\pm\frac{\alpha_sC_F}{4}(1-L_\zeta)$；同号的一对 $C^\pm_{HL}(\zeta_z)C^\pm_{HL}(\bar\zeta_z)$ 是介子准光前波函数的核，$\zeta_z=\bar\zeta_z$ 时其单圈相位为 $\pm\frac{\alpha_sC_F}{2}(1-L_\zeta)$（对应 $\mathrm{sgn}(z)=\mp1$）；$\gamma^{(1)}_C=-\frac{\alpha_sC_F}{\pi}$。(c)（两圈内容，\tagR，$R_{\mathrm{PT}}$）（$R_{\mathrm{PT}}$：在重整化微扰论中对红外正规化的部分子态严格。）在 JHEP~2023 的约定下（$\frac{\dd}{\dd\ln\mu}$，$\beta_0\equiv\beta_0^{(b)}$，$T_F=\frac12$，$\alpha_s$ 展开），% VF-YZL-7
\begin{gather}
\tilde\gamma^{(2)}_H=\frac{1188\zeta_3-99\pi^2-1108}{432\pi^2}C_FC_A-\frac{36\zeta_3+6-7\pi^2}{12\pi^2}C_F^2+\Big(\frac1{24}+\frac{10}{27\pi^2}\Big)C_Fn_f,\label{eq:yzliu-gH2}\\
\ln H^{(2)}=\frac{C_F\beta_0}{24\pi}L_\zeta^3-\frac14\Big(\Gamma^{(2)}_{\mathrm{cusp}}+\frac{C_F\beta_0}{2\pi}\Big)L_\zeta^2-\Big(\frac{C_F\beta_0}{4\pi}\Big(-2+\frac{\pi^2}{12}\Big)+\frac{\tilde\gamma^{(2)}_H}{2}\Big)L_\zeta+c_H,\label{eq:yzliu-lnH2}
\end{gather}
$c_H=(\frac{241\zeta_3}{144\pi^2}+\frac{11\pi^2}{320}-\frac{559}{1728}-\frac{971}{324\pi^2})C_FC_A+\frac{-45\zeta_3-2\pi^4+30\pi^2-30}{24\pi^2}C_F^2+\frac{36\zeta_3+51\pi^2+1312}{1296\pi^2}C_Fn_fT_F=0.0725\,C_F^2-0.0840\,C_FC_A+0.1453\,C_Fn_fT_F$，其中 $\Gamma^{(2)}_{\mathrm{cusp}}=\frac{C_F}{\pi^2}\big[C_A\big(\frac{67}{36}-\frac{\pi^2}{12}\big)-\frac59T_Fn_f\big]$ 是 $\Gamma_{\mathrm{cusp}}$ 的 $\alpha_s^2$ 系数（单圈系数为 $\frac{C_F}{\pi}$）。（本命题只涉及部分子态。）
\uses{thm:lamet-quasi-tmd,prop:yzliu-qtmd-powercount,def:sudakov-cusp,prop:lattice-renormalizability}
\end{proposition}

\begin{proofsketch}
（a）对\cref{thm:lamet-quasi-tmd}取 $\mu\frac{\dd}{\dd\mu}$：订书钉算符由端点因子乘法重整化（\cref{prop:lattice-renormalizability}），$\tilde f$ 的反常维度 $2\gamma_F$ 与 $\zeta_z$ 无关；代入 $\mu\frac{\dd\tilde K}{\dd\mu}=-2\Gamma_{\mathrm{cusp}}$、$\mu\frac{\dd\ln f^{\mathrm{TMD}}}{\dd\mu}=-\gamma_V-\Gamma_{\mathrm{cusp}}\ln\frac{\zeta}{\mu^2}$、$\mu\frac{\dd\ln S_r}{\dd\mu}=2\gamma_{HH}$，$\zeta$ 依赖相消，即得\cref{eq:yzliu-H-rge}；单圈 $\gamma_V^{(1)}=-\frac{3\alpha_sC_F}{2\pi}$、$2\gamma_F^{(1)}=\frac{3\alpha_sC_F}{2\pi}$、$\gamma_{HH}^{(1)}=-\frac{\alpha_sC_F}{\pi}$。（b）$H^{(1)}$ 是\cref{thm:lamet-quasi-tmd}(b) 的单圈内容，分解见下例；$\mu\frac{\dd H^{(1)}}{\dd\mu}=-2\partial_{L_\zeta}H^{(1)}=\frac{\alpha_sC_F}{\pi}(L_\zeta-1)$ 给出 $\Gamma^{(1)}_{\mathrm{cusp}}$ 与 $\gamma^{(1)}_C$。（c）以 $\frac{\dd L_\zeta}{\dd\ln\mu}=-2$、$\frac{\dd\alpha_s}{\dd\ln\mu}=-\beta_0\alpha_s^2$ 逐阶积分\cref{eq:yzliu-H-rge}，$\alpha_s^2$ 阶给出 $\partial_L\ln H^{(2)}=-\frac12[\Gamma^{(2)}_{\mathrm{cusp}}L+\tilde\gamma^{(2)}_H+\beta_0h_1(L)]$，$h_1=\frac{C_F}{2\pi}(-2+\frac{\pi^2}{12}-\frac{L^2}{2}+L)$，积分即\cref{eq:yzliu-lnH2}；$\tilde\gamma^{(2)}_H$ 由两圈输入\cite[式~(A.1)--(A.5)、(A.11)]{Ji:2023pba}得到，与其式~(A.13) 一致；$c_H$ 取自\cite[式~(5.11)]{Ji:2023pba}，与直接的 NNLO 计算\cite{delRio:2023pse}一致；$\tilde\gamma^{(2)}_H$ 与 $c_H$ 的 $n_f$ 部分与独立的大 $\beta_0$ 结果\cite[式~(2.41)--(2.42)]{Liu:2023onm}一致。
\end{proofsketch}

\paragraph{地位与缺口。}RGE 的结构为全阶论证 [A]（来自\cref{thm:lamet-quasi-tmd}的 RG 相容性，[A]）；(b)、(c) 在所述阶上为 [R, $R_{\mathrm{PT}}$]。勘误：文献\cite{Ji:2019ewn}式~(26) 的 $\gamma^{(2)}_C$ 有误，正确值见\cite[式~(A.12)--(A.13)]{Ji:2023pba}（附录~Q 第~14~条）。
% Verification: claim E10 (Ji:2023pba cites the earlier value without noting the discrepancy); scripts s05, s06 (W26).

\begin{example}[单圈 $H=C^+_{HL}C^-_{HL}$ 与相位]\label{ex:yzliu-hl-decomposition}
由 $L_++L_-=2L_\zeta$、$L_+^2+L_-^2=2L_\zeta^2-2\pi^2$，$(C^+C^-)^{(1)}=C^{+(1)}+C^{-(1)}=\frac{\alpha_sC_F}{8\pi}(-8-\frac{5\pi^2}{3}+4L_\zeta-2L_\zeta^2+2\pi^2)=\frac{\alpha_sC_F}{4\pi}(-4+\frac{\pi^2}{6}+2L_\zeta-L_\zeta^2)=H^{(1)}$。又 $L_+^2=L_\zeta^2-\pi^2+2\ii\pi L_\zeta$，故 $\Imag C^{+(1)}=\frac{\alpha_sC_F}{8\pi}(2\pi-2\pi L_\zeta)=\frac{\alpha_sC_F}{4}(1-L_\zeta)$，$(C^+)^2$ 的相位为 $\frac{\alpha_sC_F}{2}(1-L_\zeta)$；在 $H=C^+C^-$ 中两个相位相消，这就是准 TMD 的核为实而光前波函数的核有虚部的原因。
\uses{prop:yzliu-hard-kernel}
\end{example}
